\documentclass[11pt,a4paper]{article}
\usepackage[T1]{fontenc}
\usepackage[utf8]{inputenc}
\usepackage{lmodern}
\usepackage{amsmath,amssymb,amsthm}
\usepackage[margin=1.15in]{geometry}
\usepackage[hidelinks]{hyperref}
\emergencystretch=3em

\newtheorem{theorem}{Theorem}[section]
\newtheorem{proposition}[theorem]{Proposition}
\newtheorem{lemma}[theorem]{Lemma}
\newtheorem{corollary}[theorem]{Corollary}
\theoremstyle{definition}
\newtheorem{definition}[theorem]{Definition}
\newtheorem{remark}[theorem]{Remark}
\newtheorem*{thmA}{Theorem A}
\newtheorem*{thmB}{Theorem B}
\newtheorem*{thmC}{Theorem C}
\newtheorem*{thmD}{Theorem D}

\newcommand{\C}{\mathbb{C}}
\newcommand{\Q}{\mathbb{Q}}
\newcommand{\Z}{\mathbb{Z}}
\newcommand{\F}{\mathbb{F}}
\newcommand{\Gm}{\mathbb{G}_m}
\newcommand{\one}{\mathbf{1}}
\newcommand{\Sym}{\operatorname{Sym}}
\newcommand{\Hilb}{\operatorname{Hilb}}
\newcommand{\tr}{\operatorname{tr}}
\newcommand{\ord}{\operatorname{ord}}
\newcommand{\Spec}{\operatorname{Spec}}
\newcommand{\Stab}{\operatorname{Stab}}
\newcommand{\Hom}{\operatorname{Hom}}
\newcommand{\diag}{\operatorname{diag}}
\newcommand{\rank}{\operatorname{rank}}
\newcommand{\Tr}{\operatorname{Tr}}

\title{Counterexamples to the positivity and saturation hypotheses of GCT~VI\thanks{AI-produced. The results are due to two runs of GPT-6
Astra (OpenAI, ChatGPT Pro mode); the text was written by Claude Opus~5.5
(Anthropic) at effort high, which also supplied the steps marked as added
in Section~\ref{sec:provenance}; see that section for who proved, wrote and
checked what. No human has checked this text. To the extent that any rights
exist in it, they are dedicated to the public domain under CC0~1.0.}}
\author{}
\date{25 November 2026}

\begin{document}
\maketitle

\begin{abstract}
Part VI of Mulmuley's geometric complexity theory formulates three hypotheses
on the stretching quasi-polynomials of Kronecker coefficients, plethysm
coefficients and, more generally, multiplicities of subgroup restriction:
the saturation index (SH) and the positivity index (PH2) should be bounded by
a polynomial in the size of the input, and the generating function should
admit a positive rational form whose largest denominator exponent, the modular
index, is bounded by such a polynomial (PH3). We show that these hypotheses
fail. For every $n\ge2$ the series $\sum_{s\ge0}g((s^n),(s^n),(s^n))\,t^s$
of triple-rectangular Kronecker coefficients has a pole at a primitive root of
unity of order $2^{n-1}3^{r}$ for some $r\ge0$, so every positive rational
form, indeed every expression $P(t)/\prod_i(1-t^{a_i})^{b_i}$, has some
$a_i\ge2^{n-1}$, while the input rank is $n^2+5n$; this refutes the existence
clause of PH3 for Kronecker coefficients, also in the relaxed form that
Mulmuley gave in the appendix to the paper of Briand, Orellana and Rosas. The
proof reduces modulo three to the cyclic permutation of the three tensor
factors, whose trace is the Hilbert series of the invariants of
$\C^n\oplus\mathbb S_{(2,1)}\C^n$, and locates there a unique fixed curve of
the grading group $\mu_{2^{n-1}}$. For the plethysm coefficients of
$\Sym^{d}(\Sym^3\C^n)$ along the ray of the determinant, the Hilbert series
has simple poles at the primitive $2^{n-1}$-th roots of unity for every
$n\ge4$, which refutes PH3; a signed trace of these poles dominates all other
terms, which gives positivity and saturation indices larger than $2^{cn}$ for
every $c<1/4$ and all large $n$, refuting PH2 and SH(a). Both arguments rest
on a support theorem for cubic forms: a diagonal matrix of two-power order
fixing a semistable cubic in $n$ variables has determinant of order at most
$2^{n-1}$, with equality only for the chain cubic
$x_1^2x_2+\dots+x_{n-1}^2x_n+x_n^3$. For restriction from a rank-one torus,
which GCT~VI admits as an explicit embedding, the three hypotheses fail at a
fixed input rank.
\end{abstract}

\tableofcontents

\section{Introduction and results}\label{sec:intro}

\paragraph{The hypotheses.}
Mulmuley's GCT~VI~\cite{GCT6} proposes to decide the nonvanishing of
structure constants of representation theory, among them the Kronecker
coefficients and the plethysm coefficients, by linear programming on a
stretching quasi-polynomial, and it formulates three hypotheses that would
make this work. Let $f(n)$ be the stretching function of a structure constant
(the constant at the $n$-fold stretched input), a quasi-polynomial, and
$F(t)=\sum_{n\ge0}f(n)t^n$ its generating function. The \emph{positivity
index} $p(f)$ and the \emph{saturation index} $s(f)$ measure how far $f$ must
be shifted before all constituents have nonnegative coefficients, respectively
are positive at every positive integer whenever not identically zero; a
\emph{positive form} of $F$ is an expression with denominator
$\prod_i(1-t^{a_i})^{d_i}$ and a numerator with nonnegative integer
coefficients and constant term one, and its \emph{modular index} is
$\max_ia_i$. (The definitions are recalled in Section~\ref{sec:defs}.) The
hypotheses ask that $s(f)$ (SH(a)), $p(f)$ (PH2) and the modular index of some
positive form (PH3) be bounded by a polynomial in the rank of the input:
\cite[Hypotheses 3.3.2(a), 3.3.3, 3.3.4, p.~51]{GCT6} in general and
\cite[Hypotheses 1.6.5(a), 1.6.6, 1.6.7, pp.~24--25]{GCT6} for plethysm.
After Briand, Orellana and Rosas found counterexamples to the earlier, strict
forms of SH and PH2 for Kronecker coefficients~\cite{BOR}, Mulmuley stated
relaxed forms of SH, PH2 and PH3 for Kronecker coefficients in an appendix to
their paper~\cite[Appendix, pp.~18--20]{BOR}; the polynomial bounds above are
of that relaxed kind.

\paragraph{Results.}
Write $g(\lambda,\mu,\nu)$ for the Kronecker coefficient and put
\[
 k_n(s)=g\bigl((s^n),(s^n),(s^n)\bigr),\qquad
 K_n(t)=\sum_{s\ge0}k_n(s)t^s .
\]
This is the stretching function of the Kronecker coefficient of the triple of
columns $((1^n),(1^n),(1^n))$.

\begin{thmA}[Theorem~\ref{thm:kronecker}]
For every $n\ge2$ there is an integer $r\ge0$ such that $K_n(t)$ has a pole at
every primitive root of unity of order $2^{n-1}3^r$. Consequently every
expression $K_n(t)=P(t)/\prod_i(1-t^{a_i})^{b_i}$ with $P\in\Q[t]$ and
positive integers $a_i,b_i$ has $\max_ia_i\ge2^{n-1}$. The existence clause of
PH3 for Kronecker coefficients \cite[Hypothesis 3.3.4]{GCT6} fails, and so
does the relaxed PH3 of \cite[Appendix]{BOR}.
\end{thmA}

Let $\nu=\gcd(n,3)$, $d_0=n/\nu$ and $m=\binom{n+2}3=\dim\Sym^3\C^n$. The
plethysm coefficient of the ray $\mu=(3)$, $\lambda=(d_0)$, $\pi=((3/\nu)^n)$
of \cite[Section 1.6]{GCT6} has stretching function
\[
 f_n(s)=\dim\Hom_{GL_n}\bigl((\det)^{3s/\nu},\Sym^{d_0s}(\Sym^3\C^n)\bigr),
 \qquad H_n(t)=\sum_{s\ge0}f_n(s)t^s ,
\]
the number of linearly independent $SL_n$-invariants of degree $d_0s$ of
cubic forms in $n$ variables.

\begin{thmB}[Theorem~\ref{thm:plethysmPH3}]
For every $n\ge4$, $H_n(t)$ has a simple pole at every primitive root of unity
of order $2^{n-1}$. Every expression $H_n(t)=P(t)/\prod_i(1-t^{a_i})^{b_i}$ has
some $a_i$ divisible by $2^{n-1}$, while the parameter of
\cite[Hypothesis 1.6.7]{GCT6} is $m^2+n+1=O(n^6)$. PH3 fails for plethysm
coefficients.
\end{thmB}

\begin{thmC}[Theorem~\ref{thm:indices}]
For every real $0<c<1/4$ and all sufficiently large $n$, the positivity index
and the saturation index of $f_n$ satisfy
\[
 p(f_n)>2^{cn},\qquad s(f_n)\ge2^{cn};
\]
explicitly, $p(f_n)>2^{n/8}$ and $s(f_n)\ge2^{n/8}$ for $n\ge2\cdot10^{80}$.
PH2 and SH(a) for plethysm coefficients \cite[Hypotheses 1.6.5(a),
1.6.6]{GCT6} fail.
\end{thmC}

\begin{thmD}[Propositions~\ref{prop:torusmodular} and~\ref{prop:torusindices}]
For restriction from $\Gm$ along explicit polynomial embeddings into $GL_2$
and $GL_3$, which \cite[Section 3.4]{GCT6} admits, the input rank stays $5$,
respectively $6$, while the modular index of every positive form, respectively
the positivity and saturation indices, are unbounded. Hypotheses 3.3.2(a),
3.3.3 and 3.3.4 of \cite{GCT6} fail as stated for the subgroup restriction
problem.
\end{thmD}

\paragraph{What is not claimed.}
None of these results bears on the density clauses \cite[Hypotheses
1.6.5(b), 3.3.2(b)]{GCT6}, which ask for strict saturation of almost all
inputs; our families are sparse. We do not determine the exponent $r$ of
Theorem~A (for $n=2,3,4,5$ separate computations in the run give $r=0$; we do
not use them), and we do not bound the positivity or saturation index of the
Kronecker family. Nothing here concerns the occurrence of representations in
orbit closures, the other route of geometric complexity theory, or the
complexity of deciding the vanishing of these coefficients.

\paragraph{Method.}
The refutations of PH3 for Kronecker and for plethysm coefficients detect a pole of a Hilbert series at a root of
unity $\zeta$ of large order by the fixed locus of the grading element
$\zeta$ on the invariant-theoretic quotient: a unique fixed curve along which
the quotient is smooth forces a pole (Proposition~\ref{prop:fixedcomponentpole}),
while an isolated fixed point forces regularity (Lemma~\ref{lem:isolated}).
The fixed loci are found by a support theorem for cubic forms
(Theorem~\ref{thm:fullcubicsupport}), which singles out the chain cubic. For
Kronecker coefficients the quotient is that of three-tensors, whose fixed loci
we cannot control directly; instead the cyclic permutation $\tau$ of the three
factors acts on the invariants, its trace series agrees with $K_n$ modulo
three, and by a character identity (Theorem~\ref{thm:K023}) it is the Hilbert
series of a quotient of $\C^n\oplus\mathbb S_{(2,1)}\C^n$, where the support
theorem applies again. A transfer lemma for congruent rational series
(Theorem~\ref{thm:transfer}) carries the pole back, at the price of the unknown power $3^r$.
The torus examples of Theorem~D are elementary.

\paragraph{Earlier work.}
On 25 November 2026 we ran web searches, searched the Hexagon repository, and
went through the papers that cite \cite{GCT6} or \cite{BOR} in the OpenAlex
index, for earlier statements that PH3, PH2 or SH fails in these forms;
Section~\ref{sec:literature} lists what was searched and what was not.
Briand, Orellana and Rosas~\cite{BOR} refuted the strict forms of SH and PH2
for Kronecker coefficients; the positivity indices of their examples are at
most one, so they agree with the relaxed forms. Ikenmeyer, Mulmuley and
Walter~\cite{IMW} and Fischer and Ikenmeyer~\cite{FI} showed that deciding the
positivity of Kronecker and of plethysm coefficients is NP-hard; this concerns
the exact decision problem, not the index hypotheses. Kahle and
Micha\l{}ek~\cite{KM} showed that plethysm quasi-polynomials need not count
lattice points of polytopes of the matching dimension, which bears on the
polytope hypothesis PH1. In these searches we found no earlier refutation of
PH3, or of the relaxed SH(a) and PH2.

Section~\ref{sec:provenance} says who proved and checked what.

\section{Definitions and the hypotheses}\label{sec:defs}

\paragraph{Quasi-polynomials and indices.}
A function $f:\Z_{\ge0}\to\Z$ is a quasi-polynomial if there are a period
$l\ge1$ and polynomials $f_1,\dots,f_l$ with $f(n)=f_j(n)$ for $n\equiv j
\pmod l$. The following are \cite[Definitions 1.2.2 and 1.2.4,
pp.~13--14]{GCT6}.

\begin{definition}\label{def:indices}
The quasi-polynomial $f$ is \emph{strictly positive} if all coefficients of
all $f_j$ are nonnegative; its \emph{positivity index} $p(f)$ is the least
integer $p\ge0$ such that $n\mapsto f(n+p)$ is strictly positive. It is
\emph{strictly saturated} if, for every $j$, either $f_j$ is identically zero
or $f_j(n)>0$ for every integer $n\ge1$; its \emph{saturation index} $s(f)$ is
the least integer $s\ge0$ such that $n\mapsto f(n+s)$ is strictly saturated.
\end{definition}

\begin{definition}[{\cite[Definition 1.2.7, p.~15]{GCT6}}]\label{def:positiveform}
Let $F(t)=\sum_{n\ge0}f(n)t^n$ with $f$ a quasi-polynomial of degree $d$, not
identically zero. A \emph{positive form} of $F$ is an expression
\[
 F(t)=\frac{h_0+h_1t+\dots+h_et^e}{\prod_{i=0}^k(1-t^{a_i})^{d_i}}
\]
with $h_0=1$, all $h_i$ nonnegative integers, $a_i,d_i$ positive integers and
$\sum_id_i=d+1$. Its \emph{modular index} is $\max_ia_i$.
\end{definition}

In \cite[eq.~(1.4)]{GCT6} the numerator is printed as $h_dt^d+\dots+h_0$,
with the letter $d$ of the degree of $f$. We allow a numerator of any degree
$e$; this only enlarges the class of positive forms.

Our lower bounds on modular indices use only the denominator: they hold for
every expression $P(t)/\prod_i(1-t^{a_i})^{b_i}$ with $P\in\C[t]$, with no
condition on the numerator or on the exponents $b_i$.

\paragraph{The hypotheses.}
In \cite[Section 3.3, p.~51]{GCT6} the input $x$ of a structure constant has a
rank $\rank(x)$ and a stretching quasi-polynomial $\tilde f(x,n)$, and the
hypotheses read: (SH)(a) $s(\tilde f)=O(\mathrm{poly}(\rank(x)))$, (b) the
saturation index vanishes for almost all $x$ (Hypothesis 3.3.2); (PH2)
$p(\tilde f)=O(\mathrm{poly}(\rank(x)))$ (Hypothesis 3.3.3); (PH3) the
generating function $F(x,t)=\sum_n\tilde f(x,n)t^n$ has a positive rational
form of modular index $O(\mathrm{poly}(\rank(x)))$, and more specifically the
modular index defined in \cite[Section 4.1.1]{GCT6} is
$O(\mathrm{poly}(\rank(x)))$ (Hypothesis 3.3.4). These are stated for the
Kronecker problem, the plethysm problem and the subgroup restriction problem,
among others. For subgroup restriction along an explicit polynomial
homomorphism $\rho:H\to G$, with $H$-representation $\lambda$ and
$G$-representation $\pi$, the rank is the sum of the ranks of $H$ and $G$ and
the heights of the specified weights \cite[Section 3.4, p.~52]{GCT6}; the
Kronecker and plethysm problems are special cases.

For plethysm, \cite[Section 1.6, pp.~24--25]{GCT6} states the same hypotheses
with the parameter $h=\dim G+\operatorname{ht}(\lambda)+\operatorname{ht}(\pi)$:
SH(a) asks $s\le ch^{c'}$ (Hypothesis 1.6.5(a)), PH2 asks for a polynomial
bound on the positivity index (Hypothesis 1.6.6), and PH3 asks that the
generating function have a positive rational form with modular index bounded
by a polynomial in $h$, and more specifically that this hold for the minimal
positive form attached to an h.s.o.p.\ of lexicographically minimal degree
sequence (Hypothesis 1.6.7).

The relaxed hypotheses of \cite[Appendix, pp.~19--20]{BOR} are stated for
Kronecker coefficients $g_{\lambda,\mu}^{\pi}$ with $H=GL_n\times GL_n\subset
G=GL_{n^2}$, with the parameter $h$ the largest height of $\lambda,\mu,\pi$:
(SH)(a) $s=O(\mathrm{poly}(h))$; (PH2) $p=O(\mathrm{poly}(h))$; (PH3) the
modular index of the minimal positive form is $O(\mathrm{poly}(h))$.

\paragraph{The inputs of this paper.}
For $k_n$ we take $H=GL_n\times GL_n$, $G=GL_{n^2}$ and the columns
$\lambda=\mu=(1^n)$, $\pi=(1^n)$, stretched to $(s^n)$. Then $\rank(x)=2n+n^2+3n
=n^2+5n$ in the sense of \cite[Section 3.4]{GCT6}, and $h=n$ in the sense of
\cite[Appendix]{BOR}. For $f_n$, $G=GL(\Sym^3\C^n)=GL_m$, $\operatorname{ht}
(\lambda)=1$, $\operatorname{ht}(\pi)=n$, so $h=m^2+n+1$, and the rank of
\cite[Section 3.4]{GCT6} is $n+m+n+1$. Both are bounded by a polynomial in
$n$; so is the parameter of every hypothesis above.

